% s-090-p090-m003; original PDF 90, 91
\Aphorism{s-090-p090-m003}{94}
In der Erziehung sind wir bemüht, dem Kinde das Leben in einem Bilde darzustellen, das sich später nirgends bewährt; vor allem in einer künstlichen Einfachheit, die sich künftighin arg vermissen lässt. Um ein Beispiel zu geben: wir gewöhnen den Causalitätstrieb des Zöglings, sich von der Verknüpfung zwischen That und Erfolg, Ursache und Wirkung, Schuld und Strafe eine ganz schematische Vorstellung zu bilden, indem wir in diesen Verknüpfungen eine im Leben nie erreichte Gleichmässigkeit walten lassen. Wir befehlen dem Kinde und stellen für den Fall des Gehorsams Vortheile, für den des Ungehorsams Nachtheile in Aussicht; indem wir diese imaginären Wirkungen verwirklichen, machen wir das Kind glauben, es bestehe hier ein Connex von der Unverbrüchlichkeit eines Naturgesetzes. Um dem Kinde die Zweckmässigkeit einer Handlung zu beweisen, lassen wir den Zweck eintreten, sobald das Mittel eingetreten ist; im Leben ist zweckmässig Handeln bisweilen eine minimale Erhöhung der Chance, den Zweck eintreten zu sehen, sehr oft auch eine Erniedrigung. Kurz, wir spiegeln dem Kinde eine Welt vor, in der lauter Paare, lauter untrennbare Verbindungen A-B von Dingen existiren; in Wirklichkeit sind A und B zwei sehr complicirte Gruppen von Erscheinungen, die in einzelnen Elementen an einander grenzen. An dieser Complication und Unsicherheit der Zusammenhänge, an dem Versagen des einfachen Causalitätsbegriffes haben wir im Leben oft zu leiden; das Kind erhalten wir in einer künstlichen Simplification und würden es als schweren pädagogischen Fehler ansehen, etwa für die gleiche Handlungsweise einmal zu belohnen und das andere Mal zu strafen oder die versprochene Wirkung einer Ursache ausbleiben zu lassen; d. h. das Kind so zu behandeln wie wir Erwachsenen täglich behandelt werden.
\par

